% sections/abstract.tex
\begin{abstract}
	Penguins from the Palmer Archipelago differ in shape and size across species, and males and females differ systematically within each species. This report measures how much of a species or a body mass prediction can be recovered from physical traits alone, and how much the addition of sex improves it. Four models are compared on the \texttt{seaborn} \texttt{penguins} dataset: logistic regression and a decision tree for classification, and linear regression and a random forest for regression. Each model is trained under two feature configurations, one using only the four continuous physical traits and one adding sex. To separate the effect of sex from the effect of the particular training sample, the training data is split into two disjoint halves that share the same test set, and every configuration is evaluated under both halves. Logistic regression gains between 0.045 and 0.060 in accuracy when sex is added, and the two regressors gain between 0.048 and 0.075 in R\textsuperscript{2}. The sign of the gain is the same across both training subsets in every case, which supports the conclusion that the effect is real and not an artifact of a particular random split. The decision tree is the only model with no measurable gain, consistent with its ability to recover size-related splits directly from the physical traits. The main limitations are the small test set, a leakage in the scaler fit, and the absence of a depth constraint on the tree.
\end{abstract}
\addcontentsline{toc}{section}{Abstract}